\documentclass[./eval.tex]{subfiles}
\newcommand{\head}{\noindent Máquinas Eléctricas II |Evaluación| Apellidos y Nombres.
\rule{10cm}{0.1mm}}
\newcommand{\separador}{\noindent\rule{\paperwidth}{0.1mm}\newline}
\begin{document}
\head

 \begin{enumerate}
 \item Determinar el numero de pares de polos que posee un generador que gira a una velocidad de $n=4500rpm$ y genera una frecuencia de 300Hz.
    \item Un motor trifásico posee sus bobinas conectadas en estrella. Determinar la corriente eléctrica que absorberá de la línea si al conectarlo a una red con una tensión entre fases de $V_C=400V$ desarrolla una potencia de $P=10KW$ con un factor de potencia de $cos(\varphi)=0,8$ , Averguar la potencia reactiva $Q$ y aparente $S$ del mismo.
    \item Se conectan en estrella tres bobinas iguales a una red trifásica de $V_C=220V, f=50Hz$. Cada una de las mismas posee $R=12\Omega$ de resistencia, y $X_L=20\Omega$ de reactancia inductiva. Calcular la impedancia $Z$, corrientes simples, factor de potencia $cos(\varphi)$, potencia activa $P$ potencia reactiva $Q$ y potencia aparente  $S$.
   
    \separador
\end{enumerate}
\end{document} 